\section{Related Work}
\label{sec:related_work}
\paragraph{Self-supervised learning (SSL).} 
Self-supervised learning (SSL) has brought significant advancements in representation learning by leveraging large-scale unlabeled data~\cite{mikolov2013efficient, radford2018improving, chen2020simple, grill2020bootstrap, caron2021emerging}. In natural language processing (NLP), models like BERT~\cite{devlin2019bert} utilize masked token training strategies to learn contextual embeddings. Inspired by these successes, SSL methods have been adapted to computer vision tasks. In the vision domain, masked image modeling approaches like BEiT~\cite{bao2022beit} and MAE~\cite{he2022masked} have demonstrated remarkable performance in learning visual representations by reconstructing masked image patches. MAE introduces an asymmetric encoder-decoder architecture that reconstructs masked image patches, leading to efficient training and strong downstream performance. In this work, we adapt the MAE asymmetric encoder-decoder architecture and apply it to embeddings instead of images. Similar to how BERT captures semantic context in language by modeling token dependencies, we aim to learn the semantic context between embedding vectors for long-video understanding tasks.
\vspace{-5mm}
\paragraph{Masked autoencoders for video understanding.} VideoMAE~\cite{tong2022videomae} and VideoMAE V2~\cite{wang2023videomae} extend masked autoencoders to video data by reconstructing masked video patches, enabling efficient self-supervised learning. However, they focus on short video clips up to 32 frames, and scaling to longer videos is computationally challenging due to the quadratic complexity of attention mechanisms with respect to sequence length. A decoupling strategy is explored by
CoSeg~\cite{wang2023coseg}, yet it also operates at the \emph{frame} level and therefore, it inherits the same scaling bottleneck. In contrast to these methods, our work addresses long-video understanding by operating on sequences of embeddings and aims to learn long-range dependencies without being constrained by the number of frames. Our novelty lies in the granularity of the low-level representations, leveraging clip-level features to enable efficient training, avoiding costly frame processing.

\vspace{-5mm}
\paragraph{Long-video understanding methods.} Understanding long videos, such as movies or instructional content, is challenging due to the computational cost of processing lengthy sequences and the need to model long-range temporal dependencies. State-space models (SSMs) have been proposed to address some of these challenges. Vis4mer~\cite{islam2022long} combines a transformer encoder for extracting spatial features with a structured state-space sequence model (S4)~\cite{gu2022efficiently} decoder to efficiently aggregate representations across frames. VideoMamba~\cite{li2024videomamba} is a fully SSM-based video model that processes spatiotemporal patches directly using bidirectional Mamba blocks~\cite{Zhu2024VisionME}. While SSMs reduce computational complexity compared to standard transformers, operating directly at the frame level does not fully alleviate computational demands, limiting processing to up to 64 frames~\cite{islam2022long, li2024videomamba}. Our work focuses on enhancing long-video representations through a novel approach, specifically designed for long-form content. We evaluate our LV-MAE method on long-form video classification and regression tasks, demonstrating its effectiveness in handling extended temporal sequences. A detailed discussion of long-video language understanding methods that leverage Large Language Models (LLMs) is provided in App.~\ref{app:related}, as these approaches, while related, address different aspects of the long-video understanding challenge.
